\section{Related work}\label{sec:related}

HOPPER sits at the intersection of ordinal theory, inscription protocols and
Taproot-based custody.

\paragraph{Ordinal theory and inscriptions.}
Ordinal theory assigns a number to every sat and defines how sats move
through transactions \cite{ordinals,ord}. Inscriptions attach arbitrary
content to a sat by committing a Taproot script containing an ord envelope
and revealing it in a later transaction \cite{ord}. HOPPER adopts both ideas
but uses the inscription as a state transition, not as a collectible image.
The mint document is the frog's genesis; each crossing is a small binary
inscription on the same sat.

\paragraph{Fungible token experiments.}
BRC-20 \cite{brc20} and Runes \cite{runes} use inscriptions or OP\_RETURN
messages to encode balances and transfers. They rely on an indexer to agree
on a ledger derived from those messages. HOPPER has no fungible balance and
no global ledger: the state is attached to one sat, and the custody rules are
enforced by the script that holds it. An indexer can be wrong or unavailable
without changing which frog belongs to which sat.

\paragraph{Taproot assets and off-chain systems.}
Taproot Assets \cite{taprootassets} uses Taproot commitments to represent
assets that can move on Lightning. BitVM \cite{bitvm} explores optimistic
computation anchored in Bitcoin. HOPPER is deliberately smaller: it does not
need a bridge, a channel, or a fraud proof. The game's entire outcome
function is a four-byte block-hash modulo and a direction comparison, and the
custody rules fit in four Taproot leaves.

\paragraph{Comparison.}
The distinguishing property is the combination of three choices:
(i) identity is a sat, not a token id; (ii) every state transition is an
inscription on that sat; and (iii) the spend that publishes a transition is
also the spend that authorizes it, including the chosen direction. The first
makes the art independent of a server, the second makes the state
reconstructible from chain data, and the third makes the game self-custodial.
The cost is the one-time \textsc{prepare} transaction and the need for
ordinal-aware coin selection.

\paragraph{What HOPPER is not.}
It is not a smart-contract platform, a general-purpose asset protocol, a
privacy system, or a replacement for ord. It is a game with a small,
auditable consensus surface and a rendering path that degrades to an honest
static card.
